\section{Methods}

\subsection{Data Acquisition}

\begin{sloppypar}
We pulled all our data directly from NBA.com using the \texttt{nba\_api} Python library, which lets you query the same endpoints
that power the \texttt{stats.nba.com} website. For the 2024-25 regular season, three separate tables were downloaded at the team level.

The first was the clutch traditional stats table, which we got from the \texttt{LeagueDashTeamClutch} endpoint using the Base
measure type. This gives standard box score numbers (points, field goal percentage, turnovers, and so on) but only
for clutch situations. The second pull used the same endpoint but switched to the Four Factors measure type, which returns
Dean Oliver's four factors: effective field goal percentage, turnover rate, free throw rate, and offensive rebounding
percentage. The third dataset came from \texttt{LeagueDashTeamStats} and covered the full regular season using the Base
measure type, giving us the same box score columns but aggregated across all games, not just the close ones.

Every call used \texttt{per\_mode = ``PerGame''} and was set to the 2024-25 regular season. For the clutch-specific
calls, we also passed \texttt{clutch\_time = ``Last 5 Minutes''}, \texttt{ahead\_behind = ``Ahead or Behind''}, and
\texttt{point\_diff = 5}, which is exactly how the NBA defines a clutch situation. We added a short sleep between
requests so we did not get rate-limited. The three files were saved as \texttt{clutch\_traditional.csv},
\texttt{clutch\_four\_factors.csv}, and \texttt{regular\_traditional.csv}.
\end{sloppypar}

\subsection{Data Cleaning and Preprocessing}

\begin{sloppypar}
Once we had the raw files loaded into pandas DataFrames, we did a few rounds of cleaning before anything else.

First, we trimmed each dataset down to only the columns we actually needed. Team IDs, league IDs, and rank columns were
all dropped since they added no analytical value. The columns we kept from each file are listed in Table~\ref{tab:columns}.
After that, we ran \texttt{.str.strip()} on all the team name fields, as there were some trailing spaces that would have
broken the merge otherwise.

We also renamed every variable using a prefix system: anything from the clutch datasets got a \texttt{clutch\_} prefix,
and regular season columns got \texttt{reg\_}. So for example \texttt{W\_PCT} became either \texttt{clutch\_win\_pct} or
\texttt{reg\_win\_pct} depending on which file it came from. This made it easy to tell variables apart throughout the
rest of the analysis without having to look anything up.

Before merging, we also checked for duplicate team entries in each file. None were found. Data types looked fine across
the board; everything numeric was already float or int, so no conversion was needed.
\end{sloppypar}

\begin{table}[ht]
    \centering
    \caption{Columns kept from each source file.}
    \label{tab:columns}
    \renewcommand{\arraystretch}{1.4}
    \small
    \begin{tabular}{@{} l p{8cm} l @{}}
        \toprule
        \textbf{Dataset} & \textbf{Columns Retained} & \textbf{Prefix} \\
        \midrule
        \texttt{clutch\_traditional.csv}  &
            \texttt{TEAM\_NAME, GP, W, L, W\_PCT, PTS, FG\_PCT,} \newline
            \texttt{FG3\_PCT, FT\_PCT, REB, AST, TOV, PLUS\_MINUS} &
            \texttt{clutch\_} \\[6pt]
        \texttt{clutch\_four\_factors.csv} &
            \texttt{TEAM\_NAME, EFG\_PCT, FTA\_RATE,} \newline
            \texttt{TM\_TOV\_PCT, OREB\_PCT} &
            \texttt{clutch\_} \\[6pt]
        \texttt{regular\_traditional.csv} &
            \texttt{TEAM\_NAME, GP, W, L, W\_PCT, PTS, FG\_PCT,} \newline
            \texttt{FG3\_PCT, FT\_PCT, REB, AST, TOV, PLUS\_MINUS} &
            \texttt{reg\_} \\
        \bottomrule
    \end{tabular}
\end{table}

\subsection{Dataset Merging}

\begin{sloppypar}
We merged the three cleaned tables in two steps, both using inner joins on the team name column. First,
\texttt{clutch\_traditional} and \texttt{clutch\_four\_factors} were joined together. Then that combined table was joined
with \texttt{reg\_traditional}. Since all three files came from the same season and covered all 30 teams, no rows dropped
out during either join. The final dataset, \texttt{nba\_clutch\_merged.csv}, has one row per team, 30 rows total,
and 33 columns.
\end{sloppypar}

\subsection{Feature Engineering}

\begin{sloppypar}
We created one new variable for RQ1. To see which teams over- or under-performed in clutch situations compared to their
overall record, we subtracted regular season win percentage from clutch win percentage:
\end{sloppypar}

\begin{equation}
    \texttt{clutch\_minus\_reg\_win\_pct} \;=\; \texttt{clutch\_win\_pct} \;-\; \texttt{reg\_win\_pct}
\end{equation}

A positive value means the team won more clutch games than you would expect based on their season record. A negative
value means the opposite. This column drives the bar chart in Plot 2 of RQ1.

\subsection{Statistical Analysis}

\begin{sloppypar}
For RQ2, the main method was Pearson correlation. We used \texttt{DataFrame.corr()} in pandas to calculate the
correlation between \texttt{clutch\_win\_pct} and ten other clutch-time variables:

\begin{center}
\texttt{clutch\_fg\_pct},\ \texttt{clutch\_3p\_pct},\ \texttt{clutch\_ft\_pct},\ \texttt{clutch\_efg\_pct},\
\texttt{clutch\_tov},\ \texttt{clutch\_tov\_pct},\ \texttt{clutch\_reb},\ \texttt{clutch\_orb\_pct},\
\texttt{clutch\_ft\_rate},\ \texttt{clutch\_fta}
\end{center}

We went with Pearson because all the variables are continuous and we were interested in linear relationships. We did not
build any predictive models; the goal here was exploratory, not inferential.

For RQ3, we filtered down to just the ten teams with the highest \texttt{clutch\_win\_pct} and compared their shooting
percentages between the regular season and clutch play. That meant pairing up \texttt{reg\_fg\_pct} with
\texttt{clutch\_fg\_pct}, \texttt{reg\_3p\_pct} with \texttt{clutch\_3p\_pct}, and \texttt{reg\_ft\_pct} with
\texttt{clutch\_ft\_pct} for each of those teams.
\end{sloppypar}

